\documentclass[11pt]{article}

% Change "review" to "final" for the camera-ready, or "preprint" for a
% non-anonymous version with page numbers.
\usepackage[review]{acl}

\usepackage{times}
\usepackage{latexsym}
\usepackage[T1]{fontenc}
\usepackage[utf8]{inputenc}
\usepackage{microtype}
\usepackage{inconsolata}
\usepackage{graphicx}
\usepackage{booktabs}
\usepackage{amsmath}
\usepackage{amssymb}
\usepackage{multirow}
\usepackage{xcolor}

\graphicspath{{figures/}}

% \todo{...} marks an unverified or untraced number/claim for the critics.
\newcommand{\todo}[1]{\textcolor{red}{\textbf{[TODO: #1]}}}

% ============================================================================
% This file is maintained by the `paper-writer` subagent across loop rounds.
% Industrial-track framing: a practitioner must choose ONE multilingual
% embedding model to deploy for chemistry-patent retrieval, and must measure
% cross-lingual robustness honestly (average recall hides cross-lingual
% collapse). The paper contributes (a) two patent-grounded multilingual/CLIR
% benchmarks, (b) new robustness metrics (CLIR-MRS / MRS and the cross-lingual
% RBO / confusion / separability family), and (c) a model-by-model analysis
% yielding a concrete deployment recommendation.
% NOTE: Every number in this paper MUST trace to a file under reports/.
% ============================================================================

\title{Measuring Cross-Lingual Robustness for Chemistry-Patent Retrieval:
Patent-Grounded Multilingual Benchmarks, New Metrics, and a Deployment Decision}

\author{Anonymous Submission \\
  EMNLP Industry Track \\
  \texttt{anonymous@example.com}}

\begin{document}
\maketitle

\begin{abstract}
A chemistry-patent search team that must deploy a single multilingual
embedding model typically ranks candidates by average Recall@10---the one
number that hides the failure they will ship. Average recall is inflated by
\emph{same-language} hits; the moment a query must reach a patent in another
language (the normal case in a patent family), recall collapses, and no two
language versions of the same question return the same documents. We make this
collapse measurable. We release two patent-grounded multilingual
chemistry-retrieval benchmarks built only from human-translated patent
text---an alias-graph benchmark (132 questions, 24 ChEBI compounds, with
chemically-confusable hard negatives) and a cross-lingual (CLIR) benchmark (137
questions over a 23{,}487-document shared corpus, with Spanish as a pure
``no-home'' query language)---and a robustness-metric family that reports
cross-lingual recall, ranking consistency, confusion, and embedding
separability \emph{next to} recall. Across nine models, best cross-lingual
Recall@10 is only $0.50$, the same-language home advantage reaches $+0.55$,
five language versions of a query agree on barely a third of their top-10
(cross-lingual RBO ceiling $0.39$), and a wrong but confusable compound
out-ranks the right one on $14$--$78\%$ of queries. We show the cause is an
embedding-level separability deficit ($r(\mathrm{AUC},\mathrm{CLIR@10})=+0.96$),
so the lever is representation alignment at indexing time, not a monolingual
re-ranker at query time. \texttt{embeddinggemma} survives both benchmarks
(CLIR-MRS $0.71$), and machine-translating the \emph{question} is safe
(no significant penalty), so human-translation budget belongs on the corpus.
\end{abstract}

% ===========================================================================
\section{Introduction}
% ===========================================================================

A chemistry-patent search team that goes multilingual faces a deceptively
simple engineering decision: pick \emph{one} dense multilingual embedding model
to index a cross-jurisdictional patent corpus, and ship it. The dashboard that
guides that choice almost always shows a single number---mean Recall@10 on a
held-out query set---and the team picks the model on top. We argue that this
number is the wrong number, and that it systematically hides the failure the
team will actually deploy.

The reason is structural. In a multilingual patent corpus, the same invention
appears in several languages, and a query in one language must routinely reach
a document in another (a German chemist searching an English or Chinese family
member). Average recall, however, is dominated by the \emph{easy}
same-language case: when a same-language copy of the answer exists, every model
finds it. The hard, deployment-relevant case---retrieving a \emph{foreign}
language version---is averaged away. Our cross-lingual benchmark makes this
concrete by construction: it contains 34 Spanish queries with \emph{zero}
Spanish gold documents, a built-in ``no-home'' stress test that average recall
cannot see.

Two distinct failures hide behind that average, and both are deployment bugs.
First, \textbf{inconsistency}: ask the same compound in five languages and the
five ranked lists barely overlap---even the best model reaches a cross-lingual
rank-biased-overlap (RBO) of only $0.39$ on the alias-graph benchmark and
$0.19$ on the cross-lingual benchmark, meaning two language variants of a query
share fewer than four of their top-10 documents. Second, \textbf{confusion}: a
chemically-confusable \emph{wrong} compound out-ranks every correct patent on
$14$--$78\%$ of queries depending on the model. A search team that reports only
mean recall sees neither.

We build the instruments to measure both failures, identify the model that
survives them, and trace the cause to the embedding level. Our contributions
are:

\begin{itemize}
  \item \textbf{C1.} Two patent-grounded multilingual chemistry-retrieval
  benchmarks built only from human-translated patent text: an
  \emph{alias-graph} benchmark whose gold relevance is defined by a chemical
  ontology and whose hard negatives are confusable neighbours, and a
  \emph{cross-lingual} benchmark over a 23{,}487-document shared corpus whose
  gold is genuinely parallel human-translated patents (\S\ref{sec:benchmarks}).
  \item \textbf{C2.} A cross-lingual robustness-metric family---CLIR@k and the
  home-advantage gap, mate-retrieval, cross-lingual RBO, language-collapse, and
  embedding separability AUC, composed into a single capability$\times$robustness
  score (CLIR-MRS / MRS)---reported co-equally with recall (\S\ref{sec:metrics}).
  \item \textbf{C3.} A mechanism finding: cross-lingual failure here is an
  availability confound plus a chemistry-specific confusability trap plus an
  embedding-level separability deficit, yielding the falsifiable claim that a
  monolingual re-ranker cannot recover under-scored foreign twins
  (\S\ref{sec:analysis}).
  \item \textbf{C4.} A concrete, audited deployment recommendation with an
  operating rule: deploy \texttt{embeddinggemma}, report robustness next to
  recall, do not naively ensemble, and spend human-translation budget on the
  corpus rather than on the queries (\S\ref{sec:deployment}).
  \item \textbf{C5.} The reproducible, human-validated QAC generation and audit
  pipeline the benchmarks rest on (\S\ref{sec:benchmarks}, Appendix).
\end{itemize}

In one line: report robustness, not mean recall; and fix cross-lingual
retrieval by aligning representations at index time, not by re-ranking at query
time. Figure~\ref{fig:teaser} previews the gap that motivates the paper---every
model loses ground when no same-language document exists.

\begin{figure}[t]
  \centering
  \includegraphics[width=\linewidth]{cp_fig01_clir_leaderboard.png}
  \caption{Cross-lingual vs.\ same-language retrieval, nine models. Overall
  Recall@10 (grey) sits between the easy same-language case (MoLIR@10, green)
  and the hard cross-lingual case (CLIR@10, red). The best cross-lingual
  Recall@10 is only $0.50$ (\texttt{embeddinggemma}); average recall hides the
  collapse. Source: chem-patents key findings, Fig.~1.}
  \label{fig:teaser}
\end{figure}

% ===========================================================================
\section{Related Work}
% ===========================================================================

\paragraph{Multilingual and cross-lingual retrieval benchmarks.}
MIRACL \citep{miracl2023}, MMTEB \citep{mmteb2025}, and the TREC NeuCLIR track
\citep{neuclir2023} established large multilingual and cross-lingual retrieval
evaluations, and MTEB \citep{muennighoff2023mteb} is the de facto leaderboard
practitioners consult. These efforts report a \emph{single averaged} score per
model; our contribution is orthogonal---a \emph{retrieval-side, ranking-level}
robustness suite reported co-equally with recall, purpose-built to expose
same-language bias that an average conceals. Recent work on the contamination
and maintenance of these benchmarks \citep{mteb_contamination2025} further
motivates a fresh, content-controlled evaluation asset.

\paragraph{Cross-lingual RAG and language preference.}
BordIRlines \citep{bordirlines2024}, XRAG \citep{xrag2025}, and Linguistic
Nepotism \citep{nepotism2025} document that retrievers and RAG systems prefer
documents in particular languages. Crucially, those studies confound
\emph{content} with \emph{language}: their multilingual documents differ in
content (Wikipedia articles that are not translations of one another) or are
machine-translated English imports (translationese as a confound). Our parallel
\emph{human-translated} patents are the content-controlled asset that isolates
the language effect from the content effect---the defense for C1.

\paragraph{Calibration, conformal, and fairness in IR.}
A growing line scores the \emph{calibration} of retrieval and RAG pipelines:
TRAQ \citep{traq2023}, CONFLARE \citep{conflare2024}, and Conformal-RAG
\citep{conformalrag2025} apply conformal prediction, while fairness-via-optimal-transport
post-processing \citep{fairnessot2023,buyl2022optimal} re-balances scores. Our
suite is complementary but distinct: it scores \emph{ranking robustness}
(cross-lingual RBO, confusion, CLIR-MRS), not calibration. A reviewer familiar
with these tools should read our metrics as measuring \emph{whether the right
foreign document is ranked highly and consistently}, a question conformal
coverage does not answer.

\paragraph{Patent IR and patent-family non-equivalence.}
Patent retrieval has a long history (e.g.\ the PRIME line of NTCIR patent
experiments \citep{prime2002}). A key reason for our benchmark design is that
members of a patent family ``may differ in their claims,'' so
\texttt{publication\_number}-equivalence is an unsafe gold signal. We therefore
define gold by ontology membership (alias-graph) and by genuinely parallel
human translation (cross-lingual benchmark), never by family-number equivalence.

\paragraph{Chemistry IR and entity representations.}
SapBERT \citep{sapbert2021} aligns biomedical entity representations and PaECTER
\citep{paecter2024} learns patent-level representations, but both are largely
monolingual. We evaluate cross-lingual chemistry retrieval with ontology-graph
grounding (ChEBI, \citealp{chebi2016}) and confusable-neighbour hard negatives.

\paragraph{When to translate.}
Pre-neural CLIR studied translation-quality prediction and routing
\citep{kishida2008}; our oracle / per-language-routing discussion
(\S\ref{sec:deployment}) is the modern, embedding-level analogue of that
``when to translate'' question.

We note for completeness that calibration-flavoured cross-lingual metrics
(e.g.\ cross-lingual translation consistency or calibration-OT instruments)
are a natural \emph{future-work} horizon; we do not claim them here.

% ===========================================================================
\section{Benchmarks}
\label{sec:benchmarks}
% ===========================================================================

We release two benchmarks, both grounded in real chemistry-patent text and both
deliberately avoiding the two shortcuts that make a ``multilingual'' benchmark
misleading: machine-translated source documents (translationese) and
family-number gold (non-equivalence).

\paragraph{Shared corpus.}
Both benchmarks retrieve against one shared haystack, the
\texttt{multilingual\_GP} corpus of \textbf{23{,}487} documents drawn from
Google Patents and EPO bulk data in English, German, Spanish, French, and
Chinese, so that results are directly comparable across benchmarks.\footnote{%
Corpus row count, chem-patents key findings header.} Descriptive
corpus-construction statistics (source counts after de-duplication; IPC/A61
distribution) are reported in the supplementary system description; the
load-bearing benchmark sizes used by every result below are
\todo{trace: corpus dedup count 14{,}401 and GP/EPO/JRC coverage matrix live in
the source system slides, not yet under reports/; restrict load-bearing
numbers to the two key\_findings files until dumped}.

\paragraph{Alias-graph benchmark.}
Here a query is a multilingual ChEBI alias set for a compound (e.g.\ tin:
\emph{tin}/Sn, \emph{Zinn}, \emph{esta\~no}, \emph{\'etain}), the gold set is the
patents about that concept, and the hard negatives are documents about
chemically-confusable neighbours (siblings, parents, tautomers) taken from the
compound's alias graph. The benchmark has \textbf{132} questions over
\textbf{24} ChEBI compounds in five languages, and reports \emph{two co-equal
relevance lenses}: a \emph{concept} lens (any patent about the compound;
$\sim$109 gold per query) and a \emph{per-publication} lens (the query's own
source patent and its translations; $\sim$2.4 gold per query).\footnote{Alias
benchmark size and lens definitions, alias-graph key findings header.} The
per-publication lens is the one we read as a quality score; the concept lens is
capped by its $\sim$109 positives and must not be read as recall quality.

\paragraph{Cross-lingual (CLIR) benchmark.}
Here a query is a generated question about a source patent, and the gold set is
\emph{every language version} of that source patent. The benchmark has
\textbf{137} questions in English, German, Spanish, French, and Chinese, split
into two non-symmetric populations: \textbf{57 original} questions, asked in the
source patent's own language (so each has exactly one same-language gold), and
\textbf{80 synthetic} questions, machine-translated into another language (so
they have \emph{no same-language gold at all}---pure cross-lingual
retrieval).\footnote{Population split and Spanish property, chem-patents key
findings header and ``defining property.''} \textbf{Spanish is a pure
query-side language}: 34 Spanish queries map to zero Spanish gold documents,
giving us a controlled ``no-home'' setting in which average recall has nowhere
easy to hide.\footnote{Benchmark defining property, chem-patents key findings.}

\paragraph{Honesty by design.}
Two choices keep the benchmarks honest. (i) \emph{Source text is
human-translated only.} We never use machine-translated documents in the
corpus, so a cross-lingual miss is a retrieval failure, not a translationese
artifact. (ii) \emph{Gold is defined by ontology or by human translation, never
by \texttt{publication\_number}}, side-stepping the patent-family
non-equivalence problem (\S2). Machine translation is used only to translate
\emph{questions} for the synthetic population, and \S\ref{sec:results} shows
that this introduces no significant cross-lingual penalty.

\paragraph{Generation and validation pipeline (C5).}
The queries are produced by a reproducible per-(document, language, mode)
generation pipeline that emits three scored candidates best-first, graded on
faithfulness and answerability, across technical and semantic question modes.
A human annotator reviewed a sample of the generated items; the human-rated
quality and the LLM auto-grader's strictness relative to humans are reported in
the system description and summarized in the Appendix:
\todo{trace: human-eval numbers (e.g.\ mean 8.33/10, 97/100 reviewed, 0
rejected, auto-grader $+4.3$pp stricter) live in the source system slides, not
yet under reports/; either restrict to descriptive text or dump to reports/}.

% ===========================================================================
\section{Metrics}
\label{sec:metrics}
% ===========================================================================

Recall measures \emph{capability}; it does not measure \emph{robustness to
language}. We define a family of metrics that isolate the cross-lingual axes a
single recall number averages away. All are computed per model on both
benchmarks.

\paragraph{CLIR@k, MoLIR@k, and home advantage.}
We split Recall@k by whether the retrieved gold is in the query's own language.
\emph{MoLIR@k} (mono-lingual) restricts to same-language gold; \emph{CLIR@k}
restricts to cross-language gold. The \emph{home advantage}
$\mathrm{MoLIR}{-}\mathrm{CLIR}$ is the size of the same-language head start.
This pair is the direct instrument for ``average recall hides collapse.''

\paragraph{Directional CLIR and asymmetry.}
We compute a query-language $\rightarrow$ document-language recall matrix and
its asymmetry, because retrieval is not a symmetric similarity: the recall of
reaching language $b$ from a query in language $a$ need not equal the reverse.

\paragraph{Mate-retrieval.}
For a query whose answer also exists as a foreign \emph{twin} (a parallel
translation), we measure whether that twin is ever retrieved and at what depth:
mate-hit@k, mate-MRR, and the rank of the first foreign twin. This isolates
``can the foreign version be found at all'' from ``is it ranked first.''

\paragraph{Cross-lingual RBO.}
For a compound asked in five languages, we measure rank-biased overlap
\citep{rbo2010} between the five top-$k$ lists; RBO$=1$ means identical
rankings, $0$ means disjoint. This is our \emph{consistency} axis: does the
deployed index return the same patents regardless of query language.

\paragraph{Language collapse / over-representation.}
We measure how much a query language over-fetches documents in its own
language relative to the corpus base rate, and how often same-language
\emph{noise} out-ranks foreign \emph{gold}. This is the bias mechanism behind
the home advantage.

\paragraph{Separability AUC.}
We measure $\mathrm{AUC}(\text{score(gold)} > \text{score(non-gold)})$
separately for same-language and cross-language gold. A low cross-language AUC
means foreign gold is \emph{under-scored}, not merely mis-ordered---the
distinction that decides whether a re-ranker can help.

\paragraph{Composite: CLIR-MRS / MRS.}
We fold these into one deployable number. For the cross-lingual benchmark,
\begin{equation}
\label{eq:clirmrs}
\mathrm{CLIR\text{-}MRS} = \mathrm{cap} \times \bigl(0.5 + 0.5\cdot\mathrm{rob}\bigr),
\end{equation}
where capability $\mathrm{cap}$ aggregates accuracy, CLIR, and separability, and
robustness $\mathrm{rob}$ aggregates consistency, MT-robustness, and
language-parity; capability is the spine and robustness modulates it by
$\pm50\%$.\footnote{CLIR-MRS definition, chem-patents key findings ``Verdict.''}
For the alias-graph benchmark, the analogous MRS is the min--max-normalised mean
of five axes---accuracy, consistency, confusion-robustness, language-parity, and
separability---reported with a 95\% bootstrap CI.\footnote{MRS definition,
alias-graph key findings ``Headline numbers'' footnote.}

% ===========================================================================
\section{Experimental Setup}
\label{sec:setup}
% ===========================================================================

\paragraph{Models.}
We evaluate \textbf{nine} multilingual embedding models as drop-in dense
retrievers: \texttt{embeddinggemma}, \texttt{bge-m3} \citep{bgem3_2024},
\texttt{granite-278m}, \texttt{nomic-v2-moe}, \texttt{qwen3-0.6B},
\texttt{LaBSE} \citep{labse2022}, \texttt{SapBERT} \citep{sapbert2021},
\texttt{e5-large-instruct}, and \texttt{gte-base}.

\paragraph{Protocol.}
Both benchmarks retrieve against the single shared \texttt{multilingual\_GP}
haystack (23{,}487 documents); queries and qrels come from the respective
benchmark datasets. We report Recall@10 and the full metric family of
\S\ref{sec:metrics}. For the alias-graph benchmark we report both relevance
lenses; for the cross-lingual benchmark we report the original/synthetic split.

\paragraph{Reproducibility.}
Every figure and number is regenerated by
\texttt{reports/runs/chem\_patents/experimental\_codes/run\_all.py} and the
corresponding \texttt{reports/runs/alias\_graph/experimental\_codes/} scripts,
which also emit the executive summaries this paper cites.

% ===========================================================================
\section{Results}
\label{sec:results}
% ===========================================================================

We report the cross-lingual benchmark findings first (\S\ref{ssec:cp}), then
the alias-graph findings (\S\ref{ssec:ag}), then the two leaderboards
(\S\ref{ssec:boards}).

\subsection{The cross-lingual collapse}
\label{ssec:cp}

\paragraph{Collapse is real and large.}
Cross-lingual recall trails same-language recall in \emph{every} model
(Figure~\ref{fig:teaser}). The best CLIR@10 is $0.50$ (\texttt{embeddinggemma}),
and the home advantage $\mathrm{MoLIR}{-}\mathrm{CLIR}$ reaches $+0.55$ for the
most biased model (\texttt{e5-large-instruct}): a same-language copy of the
answer is far easier to retrieve than its foreign twin
(Figure~\ref{fig:cp_home}).

\begin{figure}[t]
  \centering
  \includegraphics[width=\linewidth]{cp_fig02_home_advantage.png}
  \caption{Same-language home advantage ($\mathrm{MoLIR}{-}\mathrm{CLIR}$) per
  model on the cross-lingual benchmark; it reaches $+0.55$.
  Source: chem-patents key findings, Fig.~2.}
  \label{fig:cp_home}
\end{figure}

\paragraph{Retrieval is anisotropic.}
The query-language $\rightarrow$ document-language recall matrix
(Figure~\ref{fig:cp_dir}) is not symmetric: the hardest direction is
en$\rightarrow$de (R@10 $=0.12$), English is the easiest \emph{target}, and the
most asymmetric pair is de$\leftrightarrow$zh (gap $+0.23$). A single ``recall''
number erases this directionality.

\begin{figure}[t]
  \centering
  \includegraphics[width=\linewidth]{cp_fig03_directional_clir_matrix.png}
  \caption{Query$\rightarrow$document language recall matrix (best model);
  retrieval direction matters. Source: chem-patents key findings, Fig.~3.}
  \label{fig:cp_dir}
\end{figure}

\paragraph{Machine-translating the question is safe.}
Controlling for the source patent, the paired human$-$MT difference in
cross-lingual reach is $-0.044$ ($p=0.13$): statistically
\emph{insignificant} (Figure~\ref{fig:cp_mt}). We read this as a \emph{null}
result---MT of the query does not significantly hurt and we do not claim it
helps---which validates the project's ``MT-is-fine-for-the-question''
assumption.

\begin{figure}[t]
  \centering
  \includegraphics[width=\linewidth]{cp_fig05_mt_penalty.png}
  \caption{Paired human vs.\ machine-translated queries; the cross-lingual
  difference is $-0.044$ ($p=0.13$), not significant.
  Source: chem-patents key findings, Fig.~5.}
  \label{fig:cp_mt}
\end{figure}

\paragraph{Foreign twins are buried.}
Pooled mate-hit@10 is $0.38$, and on $15\%$ of (query, model) pairs the foreign
twin never surfaces even in the top-1000 (Figure~\ref{fig:cp_mate}). The best
twin-finder, \texttt{embeddinggemma}, reaches its first foreign twin at a median
rank of $5$ (Figure~\ref{fig:cp_rank})---most models do far worse.

\begin{figure}[t]
  \centering
  \includegraphics[width=\linewidth]{cp_fig06_mate_retrieval.png}
  \caption{Mate-retrieval: mate-hit@10/100 and mate-MRR. Pooled mate-hit@10 is
  $0.38$. Source: chem-patents key findings, Fig.~6.}
  \label{fig:cp_mate}
\end{figure}

\begin{figure}[t]
  \centering
  \includegraphics[width=\linewidth]{cp_fig07_first_foreign_rank.png}
  \caption{Depth to the first foreign twin; $15\%$ of pairs never reach one in
  the top-1000. Source: chem-patents key findings, Fig.~7.}
  \label{fig:cp_rank}
\end{figure}

\subsection{Inconsistency and confusion (alias-graph)}
\label{ssec:ag}

\paragraph{Same question, different rankings.}
Asking the same compound in five languages does \emph{not} return the same
patents. Even the best model reaches a cross-lingual RBO of only $0.39$
(Figure~\ref{fig:ag_rbo}); two language variants of a query share fewer than
four of their top-10 documents, and Chinese is the consistent odd-one-out. The
embedding space is far from language-agnostic.

\begin{figure}[t]
  \centering
  \includegraphics[width=\linewidth]{ag_fig1_cross_lingual_rbo.png}
  \caption{Same compound, five languages: cross-lingual RBO. The ceiling is
  $0.39$ (\texttt{embeddinggemma}). Source: alias-graph key findings, Fig.~1.}
  \label{fig:ag_rbo}
\end{figure}

\paragraph{Confusable look-alikes out-rank the gold.}
A chemically-confusable wrong compound out-ranks every gold patent on
$14$--$78\%$ of queries (per-publication lens) depending on the model, worst for
German and Chinese, and driven almost entirely by \emph{sibling} compounds
rather than broader parent classes (Figure~\ref{fig:ag_conf}). The repeat
offenders are a small set of universal attractors---\emph{polypeptide},
\emph{methyl}, \emph{ethene}, \emph{hydroxide}, \emph{dioxygen}
(Figure~\ref{fig:ag_attr}).

\begin{figure}[t]
  \centering
  \includegraphics[width=\linewidth]{ag_fig2_confusion_both_lenses.png}
  \caption{Confusion rate (a look-alike beats all gold), model $\times$ language,
  both lenses; $14$--$78\%$ on the publication lens.
  Source: alias-graph key findings, Fig.~2.}
  \label{fig:ag_conf}
\end{figure}

\begin{figure}[t]
  \centering
  \includegraphics[width=\linewidth]{ag_fig5_universal_attractors.png}
  \caption{Universal attractors: look-alikes that most often steal the top rank.
  Source: alias-graph key findings, Fig.~5.}
  \label{fig:ag_attr}
\end{figure}

\subsection{Leaderboards}
\label{ssec:boards}

Table~\ref{tab:cp_board} ranks the nine models by CLIR-MRS on the cross-lingual
benchmark and Table~\ref{tab:ag_board} by MRS on the alias-graph benchmark.
\texttt{embeddinggemma} tops both. The radar views
(Figures~\ref{fig:cp_radar},~\ref{fig:ag_radar}) show \emph{where} each top
model wins across the axes.

\begin{table}[t]
  \centering
  \small
  \begin{tabular}{lrrrrr}
    \toprule
    model & R@10 & CLIR@10 & home & sep & CLIR-MRS \\
    \midrule
    \texttt{embeddinggemma} & 0.544 & 0.502 & $+0.22$ & 0.90 & \textbf{0.71} \\
    \texttt{bge-m3}         & 0.476 & 0.437 & $+0.23$ & 0.87 & 0.70 \\
    \texttt{granite-278m}   & 0.359 & 0.329 & $+0.17$ & 0.86 & 0.67 \\
    \texttt{nomic-v2-moe}   & 0.443 & 0.381 & $+0.34$ & 0.86 & 0.67 \\
    \texttt{qwen3-0.6B}     & 0.468 & 0.433 & $+0.19$ & 0.86 & 0.58 \\
    \texttt{LaBSE}          & 0.277 & 0.238 & $+0.21$ & 0.81 & 0.40 \\
    \texttt{SapBERT}        & 0.212 & 0.179 & $+0.18$ & 0.79 & 0.40 \\
    \texttt{e5-large-instruct} & 0.178 & 0.077 & $+0.55$ & 0.71 & 0.26 \\
    \texttt{gte-base}       & 0.004 & 0.000 & $+0.02$ & 0.49 & 0.00 \\
    \bottomrule
  \end{tabular}
  \caption{Cross-lingual benchmark leaderboard. ``home'' is the home advantage
  $\mathrm{MoLIR}{-}\mathrm{CLIR}$; ``sep'' is separability AUC.
  \texttt{embeddinggemma} wins (CLIR-MRS $0.71$, CI $[0.67,0.77]$).
  Source: chem-patents key findings, Verdict table / \texttt{headline\_numbers.csv}.}
  \label{tab:cp_board}
\end{table}

\begin{table}[t]
  \centering
  \small
  \begin{tabular}{lrrrr}
    \toprule
    model & pubR@10 & RBO & conf & MRS \\
    \midrule
    \texttt{embeddinggemma} & 0.67  & 0.387 & 0.144 & \textbf{0.991} \\
    \texttt{qwen3-0.6B}     & 0.641 & 0.349 & 0.227 & 0.902 \\
    \texttt{bge-m3}         & 0.604 & 0.288 & 0.237 & 0.813 \\
    \texttt{LaBSE}          & 0.511 & 0.314 & 0.295 & 0.779 \\
    \texttt{granite-278m}   & 0.529 & 0.289 & 0.318 & 0.757 \\
    \texttt{nomic-v2-moe}   & 0.580 & 0.239 & 0.215 & 0.727 \\
    \texttt{SapBERT}        & 0.352 & 0.142 & 0.424 & 0.494 \\
    \texttt{e5-large-instruct} & 0.237 & 0.021 & 0.481 & 0.239 \\
    \texttt{gte-base}       & 0.047 & 0.003 & 0.782 & 0.000 \\
    \bottomrule
  \end{tabular}
  \caption{Alias-graph benchmark leaderboard. ``pubR@10'' is per-publication
  Recall@10; ``conf'' is publication-lens confusion rate.
  \texttt{embeddinggemma} wins (MRS $0.991$, CI $[0.86,1.00]$).
  Source: alias-graph key findings, headline table / \texttt{headline\_numbers.csv}.}
  \label{tab:ag_board}
\end{table}

\begin{figure}[t]
  \centering
  \includegraphics[width=\linewidth]{cp_fig14_robustness_radar.png}
  \caption{Robustness profiles, cross-lingual benchmark.
  Source: chem-patents key findings, Fig.~14.}
  \label{fig:cp_radar}
\end{figure}

\begin{figure}[t]
  \centering
  \includegraphics[width=\linewidth]{ag_fig10_robustness_radar.png}
  \caption{Robustness profiles, alias-graph benchmark.
  Source: alias-graph key findings, Fig.~10.}
  \label{fig:ag_radar}
\end{figure}

% ===========================================================================
\section{Analysis}
\label{sec:analysis}
% ===========================================================================

Why does cross-lingual retrieval collapse here, and what does the failure tell
us about the fix?

\paragraph{The lead is an availability confound.}
English's apparent advantage is mostly \emph{where the gold lives}, not encoder
skill. Each language retrieves its own-language gold far better
($0.63$--$0.82$) than foreign-language gold ($0.35$--$0.47$), and because
English and French source patents dominate, $42\%$ of an English query's gold is
reachable in-English versus only $8$--$10\%$ for German, Spanish, and
Chinese.\footnote{Availability shares, alias-graph key findings ``mechanism.''}
The collapse is visible as language over-fetch: low-resource query languages
over-fetch their own language up to $49\times$ the corpus base rate, and
same-language \emph{noise} out-ranks the gold on $60\%$ of queries
(Figures~\ref{fig:cp_collapse},~\ref{fig:cp_distractor}).

\begin{figure}[t]
  \centering
  \includegraphics[width=\linewidth]{cp_fig09_language_collapse.png}
  \caption{Same-language over-representation by query language; up to $49\times$
  the base rate. Source: chem-patents key findings, Fig.~9.}
  \label{fig:cp_collapse}
\end{figure}

\begin{figure}[t]
  \centering
  \includegraphics[width=\linewidth]{cp_fig10_distractor_language.png}
  \caption{Language of the documents that bury the gold; same-language noise
  out-ranks the gold on $60\%$ of queries.
  Source: chem-patents key findings, Fig.~10.}
  \label{fig:cp_distractor}
\end{figure}

\paragraph{Bias drives inconsistency.}
The more a model leans on its home language, the less consistent it is across
languages: $r(\text{home advantage},\text{RBO})=-0.85$ on the cross-lingual
benchmark (Figure~\ref{fig:cp_consist}) and $r(\text{same-language
share},\text{RBO})=-0.87$ on the alias-graph benchmark
(Figure~\ref{fig:ag_bias}). The two failures are the same failure seen from two
angles.

\begin{figure}[t]
  \centering
  \includegraphics[width=\linewidth]{cp_fig08_consistency_vs_bias.png}
  \caption{RBO consistency vs.\ home advantage; $r=-0.85$.
  Source: chem-patents key findings, Fig.~8.}
  \label{fig:cp_consist}
\end{figure}

\begin{figure}[t]
  \centering
  \includegraphics[width=\linewidth]{ag_fig4_bias_drives_inconsistency.png}
  \caption{Language bias $\Rightarrow$ cross-lingual inconsistency; $r=-0.87$.
  Source: alias-graph key findings, Fig.~4.}
  \label{fig:ag_bias}
\end{figure}

\paragraph{Structure-style questions are the trap.}
\emph{How} a query is phrased decides whether the model falls into a confuser.
\emph{Structure}-style questions are the trap (Recall@10 $0.26$, confusion
$51\%$) because structure descriptions are exactly what sibling compounds
share, whereas \emph{role} questions are far safer (Recall@10 $0.60$, confusion
$25\%$); a language-independent \emph{formula token} (e.g.\ H$_2$S, CO$_2$)
measurably helps ($p<0.01$), even in Chinese
(Figure~\ref{fig:ag_qtype}).

\begin{figure}[t]
  \centering
  \includegraphics[width=\linewidth]{ag_fig6_question_type_effect.png}
  \caption{Structure questions are the trap (R@10 $0.26$, confusion $51\%$) vs.\
  role questions ($0.60$, $25\%$). Source: alias-graph key findings, Fig.~6.}
  \label{fig:ag_qtype}
\end{figure}

\paragraph{It is a separability deficit, so re-ranking cannot fix it.}
Behind the confusion, the cosine score barely separates gold from sibling:
$\mathrm{AUC}(\text{gold}>\text{look-alike})=0.55$ for confused queries versus
$0.70$ otherwise (Figure~\ref{fig:ag_sep}). On the cross-lingual benchmark the
same story holds across models: $r(\text{cross-language AUC},\text{CLIR@10})
=+0.96$ (Figure~\ref{fig:cp_sep})---foreign gold is \emph{under-scored}, not
merely mis-ordered. This is the crux: a monolingual re-ranker reads a ranked
list it cannot repair, because the right foreign document was never scored high
enough to appear. The lever is at the embedding level.

\begin{figure}[t]
  \centering
  \includegraphics[width=\linewidth]{cp_fig11_separability.png}
  \caption{Separability AUC (gold$>$non-gold), same vs.\ cross language;
  $r(\mathrm{AUC},\mathrm{CLIR@10})=+0.96$.
  Source: chem-patents key findings, Fig.~11.}
  \label{fig:cp_sep}
\end{figure}

\begin{figure}[t]
  \centering
  \includegraphics[width=\linewidth]{ag_fig8_confusion_is_separability.png}
  \caption{Confusion is a separability collapse: AUC $0.55$ confused vs.\ $0.70$
  otherwise. Source: alias-graph key findings, Fig.~8.}
  \label{fig:ag_sep}
\end{figure}

% ===========================================================================
\section{Deployment Recommendation}
\label{sec:deployment}
% ===========================================================================

The metrics yield a single, actionable decision for a search team.

\paragraph{Deploy \texttt{embeddinggemma}.}
It wins both composites (CLIR-MRS $0.71$, CI $[0.67,0.77]$; alias MRS $0.991$,
CI $[0.86,1.00]$), tops cross-lingual Recall@10 ($0.50$), and is the best
twin-finder (median first-foreign rank $5$). Mean recall alone would already
rank it first here, but---critically---it is robustness, not recall, that keeps
it first on the alias-graph benchmark, where mean recall under-rates
language-balanced encoders.

\paragraph{Report robustness next to recall.}
We recommend reporting CLIR@10, language-parity, and MRS \emph{alongside}
recall, never average recall alone. The pure ``no-home'' Spanish queries are the
reason: a dashboard that averages over them will look healthy while the deployed
system silently fails on cross-jurisdiction search.

\paragraph{Do not reflexively ensemble.}
There \emph{is} real headroom---the oracle (any of nine models) reaches CLIR@10
$0.61$ on the cross-lingual benchmark (Figure~\ref{fig:cp_ens}) and $88\%$
versus the best single model's $76\%$ on the alias-graph benchmark, with Chinese
having the largest recoverable headroom (Figure~\ref{fig:ag_ens}). But this
headroom is \emph{not free}: untuned reciprocal-rank fusion \citep{rrf2009} did
\emph{not} beat the best single model. The two facts coexist---naive fusion
loses, yet a score-aware learned combiner or per-language routing (for the
homeless Spanish and the low-resource Chinese) could win. A residual
$\sim$12\% ``universal-blind'' core (the methyl/sulfide traps) needs
chemistry-aware help, not more encoders.

\begin{figure}[t]
  \centering
  \includegraphics[width=\linewidth]{cp_fig12_ensemble_headroom.png}
  \caption{Best-single vs.\ fusion vs.\ router vs.\ oracle (cross-lingual);
  oracle CLIR@10 $0.61$, but untuned RRF does not beat the best single model.
  Source: chem-patents key findings, Fig.~12.}
  \label{fig:cp_ens}
\end{figure}

\begin{figure}[t]
  \centering
  \includegraphics[width=\linewidth]{ag_fig7_ensemble_headroom.png}
  \caption{Oracle vs.\ best single model, per language (alias-graph); $88\%$ vs.\
  $76\%$, Chinese has the largest headroom.
  Source: alias-graph key findings, Fig.~7.}
  \label{fig:ag_ens}
\end{figure}

\paragraph{Align, do not re-rank.}
Because foreign gold is under-scored at the embedding level
(\S\ref{sec:analysis}), the durable fix is representation alignment at indexing
time. A monolingual re-ranker at query time operates on a candidate list that
already excludes the right foreign document.

\paragraph{Budget rule.}
Finally, the MT-of-question null result ($-0.044$, $p=0.13$) tells us where to
spend a finite human-translation budget: \emph{machine-translate the query, but
human-translate the corpus}. The expensive resource buys the most where it
matters---the documents being searched, not the questions.

% ===========================================================================
\section*{Limitations}
% ===========================================================================

\paragraph{Scale.}
The benchmarks comprise 132 (alias-graph) and 137 (cross-lingual) questions over
24 compounds; some language-pair cells are thin, so the directional CLIR matrix
should be read as indicative rather than definitive in the sparsest cells.

\paragraph{Domain transfer.}
All results are chemistry-patent-specific. We frame that specificity as the
contribution---a real deployment niche underserved by general benchmarks---but
we do not yet provide a general-domain companion measurement:
\todo{trace: no MMTEB / MIRACL / NeuCLIR transfer number computed yet;
flagged in needs\_eval as a candidate}.

\paragraph{Judge dependence.}
The generated queries and the auto-grader rely on an LLM judge. We provide human
validation of a sample (Appendix), but claim-level patent equivalence is not
lay-annotatable, and a larger expert re-annotation remains future work.

\paragraph{Language coverage.}
We cover five languages (en/de/es/fr/zh) with a single non-Latin script
(Chinese); low-resource scripts are out of scope here.

\paragraph{Out of scope.}
Calibration-style cross-lingual instruments and a full claim-equivalence audit
are deliberately left as future work; we do not claim them as results.

% ===========================================================================
\section{Conclusion}
% ===========================================================================

For a chemistry-patent search team deploying one multilingual embedding model,
average Recall@10 is the wrong dashboard: it is inflated by same-language hits
and hides a cross-lingual collapse (best CLIR@10 $0.50$; home advantage up to
$+0.55$; five language versions of a query agreeing on barely a third of their
top-10; a confusable compound beating the gold on $14$--$78\%$ of queries). Our
two patent-grounded benchmarks and our CLIR robustness suite make that collapse
measurable, name the survivor (\texttt{embeddinggemma}), and trace the cause to
an embedding-level separability deficit---so the fix is alignment at index time,
not re-ranking at query time, and the budget rule is to machine-translate the
query and human-translate the corpus. Reporting robustness next to recall, and
treating cross-lingual retrieval as a representation problem, are the two habits
we recommend any multilingual patent-search team adopt.

\bibliography{custom}

\appendix
\section{Appendix}

\paragraph{Reproducibility.}
The full ten-round derivations behind the headline numbers, the per-language and
per-lens tables, and all figures are produced by the scripts under
\texttt{reports/runs/\{chem\_patents,alias\_graph\}/experimental\_codes/} and
summarised in the two \texttt{key\_findings/EXECUTIVE\_SUMMARY.md} files this
paper cites.

\paragraph{Generation-pipeline validation (C5).}
A human annotator reviewed a sample of generated queries; the quantitative
human-rating summary and the LLM auto-grader's strictness relative to humans are
reported in the system description and will be folded in here once dumped to
\texttt{reports/}:
\todo{trace: human-eval numbers (mean rating, items reviewed, rejections,
auto-grader strictness) currently in source slides only, not under reports/}.

\end{document}
